质量为 $m>0$ 的非相对论粒子处在对称有限方势阱
\[
V(x)=\begin{cases}-V_0,&|x|<a,\\0,&|x|\ge a,\end{cases}\qquad V_0>0,\ a>0 .
\]
只研究平方可积束缚态 $-V_0<E<0$，忽略自旋与其他作用。
(1) 利用中心对称宇称和 $x=\pm a$ 处波函数及导数连续性，建立偶态、奇态的能级超越方程；(2) 用无量纲阱强度 $\lambda=a\sqrt{2mV_0}/\hbar$ 求束缚态数目，并说明阈值处为何不把 $E=0$ 计为束缚态；(3) 给出偶、奇态波函数的归一化常数表达式，证明能级均简单而没有奇偶简并；(4) 在 $V_0\to\infty$ 的固定低能级序号极限及 $V_0\to0^+$ 的浅阱极限求能量的主导近似。
